\documentclass[11pt,a4paper]{article}
\usepackage[margin=23mm]{geometry}
\usepackage{amsmath,amssymb,graphicx,booktabs,array,longtable,url}
\usepackage[hidelinks]{hyperref}
\usepackage{fancyhdr}
\newcommand{\norm}[1]{\left\lVert#1\right\rVert}
\input{generated/metrics}
\pagestyle{fancy}\fancyhf{}
\fancyhead[L]{\small QCeNN-MI: emulator evidence and hardware validation plan}
\fancyfoot[L]{\small Technical report; physical execution pending}
\fancyfoot[R]{\thepage}
\setlength{\headheight}{15pt}
\setlength{\emergencystretch}{2em}
\setlength{\parindent}{0pt}\setlength{\parskip}{5pt}
\begin{document}
\begin{titlepage}
\vspace*{15mm}
{\Large\bfseries EuroHPC Quantum Access Pilot\par}
\vspace{8mm}
{\huge\bfseries QCeNN-MI\par}
\vspace{5mm}
{\LARGE Emulator Benchmark Evidence and\par Physical-Hardware Validation Plan\par}
\vspace{12mm}
\begin{tabular}{@{}p{43mm}p{112mm}@{}}
Application & \textbf{\ApplicationID}\\[3mm]
Project team & Vahid Tavakkoli; Kabeh Mohsenzadegan; Kyandoghere Kyamakya\\[3mm]
Report date & \ReportDate\\[3mm]
Software source & \url{https://github.com/vtavakkoli/QCeNN-MI}\\[3mm]
Source revision & \texttt{\SourceCommit} (full provenance in the artifact)\\[3mm]
Evidence status & Simulator benchmarks only; no physical QPU job claimed\\
\end{tabular}
\vfill
\textbf{Purpose.} Provide a reproducible technical-feasibility record for
assessment of a subsequent small-instance physical quantum experiment.
The report supports review; it does not assert that EuroHPC has granted
physical access or that the reported workloads ran on allocated EuroHPC nodes.
\vspace{8mm}
\end{titlepage}

\section*{Executive assessment}
QCeNN-MI now has a unified, containerized emulator workflow covering its
small hybrid matrix-inversion milestones and a separate coherent inverse-action
baseline. The complete recorded software suite passed \TestCount{} unit tests
and its numerical campaign gates. The compact coherent study audits
\PrecisionCaseCount{} matrices of order 1000, including SPD, nonsymmetric
and indefinite inputs with condition numbers up to 1000. Its largest full
operator residual is \WorstOperatorResidual{} and its largest relative
Frobenius inverse error is \WorstInverseError{}.

These high-precision results use ideal structured gate-statevector simulation
and classical dense spectral/SVD encoding. They are not native noisy
order-1000 execution. Native-gate ideal checks and synthetic-noise experiments
cover orders 2, 4, 10 and 100. The dense order-100 noisy verifier fails the
project quality gate, so the evidence supports a staged small-instance
hardware study rather than an assertion of hardware-ready large inversion.

\textbf{Requested next technical step:} consideration of a bounded
order-2/order-4 physical experiment after completion of hosting-system
reproduction, device-specific compilation and calibration-derived emulation.
Order 10 is a conditional follow-on. No dense order-1000 QPU submission is
proposed as the initial physical test.

\section{Allocation context and evidence provenance}
The applicant-provided emulator allocation information is summarized below.
It is separate from resource use measured by this software run.

\begin{tabular}{@{}p{67mm}p{90mm}@{}}\toprule
Field & Applicant-provided value\\\midrule
Application identifier & \ApplicationID\\
Emulator allocation & \AllocationNodeHours{} node-hours\\
Storage allocation & \AllocationStorageGB{} GB\\
Allocation duration & \AllocationDays{} days\\
Information source & \AllocationSource\\
Hosting system / scheduler account & Not present in the evidence supplied to this run\\
Physical QPU allocation / job IDs & None recorded\\\bottomrule
\end{tabular}

The public Quantum Access Pilot describes technical-feasibility, testing,
performance-data collection and workflow development as purposes of access
\cite{eurohpc2026}. This report is organized around those purposes. It is
not an official EuroHPC report template and does not replace instructions
from the assigned hosting entity or evaluator.

The measured run has execution label \texttt{\RunContext}, Python
\PythonVersion{}, NumPy \VersionNumpy{}, SciPy \VersionScipy{}, Qiskit
\VersionQiskit{} and Qiskit Aer \VersionQiskitaer{}. The full provenance
manifest records the source revision, source-tree hash, environment,
thread controls, commands and return codes. Its explicit fields make no
claim of execution on allocated EuroHPC nodes or on a physical QPU.
Raw timestamps use the execution environment's clock, which is not independently
synchronized here. The report date is separate from those process timestamps;
only a hosting-system scheduler record can attest an allocated job's timing.

\textbf{Source-tree SHA-256:}
\begin{quote}\small\ttfamily\raggedright
\SourceFingerprint
\end{quote}

\section{Scientific claim and theoretical basis}
The coherent baseline targets $\lvert b\rangle\mapsto
\lvert M^{-1}b\rangle$ after ancilla postselection, for a real nonsingular
square matrix. For such matrices, $MM^T\succ0$ is equivalent to invertibility.
It does not bound the condition number or guarantee fixed-resource
performance. The implementation does not solve normal equations.

Let $M=U\Sigma V^T$ and let $0<a\leq\sigma_j\leq b$ bound its singular values.
Using $c=(a+b)/2$, $h=(b-a)/2$ and
$r=(\sqrt b-\sqrt a)/(\sqrt b+\sqrt a)$, the degree-$d$ reciprocal approximation is
\[
 p_d(\sigma)=\frac{1}{\sqrt{ab}}
 \left[1+2\sum_{k=1}^d(-r)^kT_k\!\left(\frac{\sigma-c}{h}\right)\right].
\]
The accompanying IEEE-style manuscript gives the full derivation. Its main
exact-arithmetic residual guarantee is
\[
 X_d=Vp_d(\Sigma)U^T,\qquad
 \norm{MX_d-I}_2\leq\delta_d
 =2\sqrt\kappa\frac{r^{d+1}}{1-r},\qquad \kappa=b/a.
\]
It also proves relative per-input solution error at most $\delta_d$,
normalized-state error at most $2\delta_d$, and, when $\delta_d<1$,
success probability at least $(1-\delta_d)^2/\kappa^2$.

The construction specializes established reciprocal-series/LCU ideas
\cite{harrow2009,childs2017,gilyen2019}; it is not a claim of a new general
quantum linear-system algorithm. Classical dense decomposition, coefficient
and gate synthesis, and input loading are charged explicitly. The coherent
solver receives no fitted solution or classical residual correction. A
classical reference is used only after simulation for verification.

An order-1000 condition-number-10000 SPD probe needs degree 1842 at the
$10^{-12}$ truncation target and is rejected by the default degree budget
1023. This concrete limit is retained in the evidence. Raising the budget
does not remove floating-point, gate-depth or sampling constraints.

\section{Milestone-to-evidence mapping}
{\small
\begin{longtable}{@{}p{27mm}p{72mm}p{52mm}@{}}\toprule
Milestone & Measured implementation / result & Scope and outstanding evidence\\\midrule
A1 & One-qubit real inverse-column representation; ideal residual and
finite-shot X/Z reconstruction. & Small hybrid representation test; not
standalone coherent inversion.\\[3mm]
A2 & Shared coupled quantum ansatz; depth sweep 1, 2 and 4; ideal,
finite-shot and synthetic-noise tomography. & Optimized classically;
not a measured physical local quantum network.\\[3mm]
A3 & Stateful warm-start tracking and a dynamic order-2 study. & Predictor
and safeguarding are classical; physical time-varying tracking pending.\\[3mm]
A4 & Safeguarded classical polynomial correction and small dense quantum
polynomial dilations. & Classical and quantum contributions labeled;
not efficient generic polynomial block encoding.\\[3mm]
Coherent baseline & Neumann-LCU and compact Chebyshev-LCU inverse action;
small native cross-checks and ideal structured order-1000 audit. & No
classical solution fitting in the coherent core; dense spectral encoding remains.\\[3mm]
Hardware preparation & Native gate counts, restored layouts, QPY exports,
raw counts, echo confidence intervals and synthetic-noise outcomes. & No
calibration-derived noise or physical backend job is claimed.\\\bottomrule
\end{longtable}
}

The complete default depth sweep above is 1/2/4 layers. The repository can
request longer sweeps, including 8/16/32/64/150, but this report does not
claim those longer coupled-optimization experiments were completed. They
must be run and archived separately before being described as evidence.

\subsection{Small-matrix emulator convergence and noisy depth sweep}
The following rows show uncorrected ideal, finite-shot and noisy residuals.
Classical correction is retained separately in the JSON and is not used
to relabel the noisy quantum output as an accurate coherent solution.

\begin{center}\small
\begin{tabular}{llrrr}\toprule
Case & Layers & Ideal residual & Shot residual & Noisy residual\\\midrule
\input generated/a2_rows.tex 
\bottomrule\end{tabular}
\end{center}

\section{High-precision coherent order-1000 evidence}
The precision experiment uses four earlier well-conditioned families plus
nine dense controlled-spectrum cases. The polynomial target is $10^{-12}$;
the independent numerical gate is $10^{-10}$. A complete simulated operator
audit is required, in addition to three individual inputs. The audit
reconstructs the success block from simulator amplitudes and measures
$\norm{MX-I}_2$, normalized Frobenius residual and relative inverse error.
This is stronger than a few input checks, but is not physical tomography.

\begin{center}\scriptsize
\begin{tabular}{lrrrrrr}\toprule
Family & $\kappa$ & $d$ & Qubits & Operator residual & Inverse error $F$ & Min success\\\midrule
\input generated/precision_rows.tex 
\bottomrule\end{tabular}
\end{center}

The register counts comprise ten data qubits, one walk flag and binary
selectors. They exclude error-correction overhead. The order-1000
condition-number-1000 cases use 21 raw qubits, but their worst-input
success probability is approximately $10^{-6}$. With independent
postselection, that means roughly a million attempts per accepted
output for that input. Hardware utility therefore cannot be inferred
from fitting below a 42-qubit raw capacity.

\begin{figure}[htbp]\centering
\includegraphics[width=0.94\textwidth]{figures/convergence.pdf}
\caption{Theory-bound regression from simulated controlled rotations.
The experiment checks \TheoryCaseCount{} order/degree/condition/family cases.
Low-degree inaccurate cases remain visible.}
\end{figure}

\section{Native circuit preparation and negative noisy evidence}
The generic hardware-preparation campaign uses an active-width bidirectional
line and the native set $R_z$, SX, X, CX. It restores the compiler's logical
layout before calculating ideal residuals. It must be regenerated against
the actual provider target before a QPU submission.

\begin{center}\small
\begin{tabular}{rrrrr}\toprule
$N$ & Total qubits & Native CX & Native depth & Ideal residual\\\midrule
\input generated/hardware_rows.tex 
\bottomrule\end{tabular}
\end{center}

The echo verifier appends an independently computed reference projection
after inverse action and counts reference returns conditional on ancilla
success. Its gate and readout noise are part of the measured fraction.
It is not a direct state-fidelity estimate or a signed-inverse reconstruction.
The quality rule requires at least 20 accepted shots and a 95\% Wilson
lower confidence bound of at least 0.9 \cite{wilson1927}.

\begin{center}\small
\begin{tabular}{rlrrrll}\toprule
$N$ & Profile & Shots & Accepted & Return & 95\% interval & Gate\\\midrule
\input generated/noise_rows.tex 
\bottomrule\end{tabular}
\end{center}

The optimistic/stress profiles use SX/X errors $10^{-5}/10^{-4}$, CX errors
$10^{-4}/10^{-3}$ and readout flips $0.005/0.02$. They are synthetic,
not measured device calibrations. Small cases use exact density-matrix
noise channels with 32768 measurement shots. The order-100 optimistic
case uses 64 noisy statevector trajectories; no corresponding stress
case or high-precision confidence claim is implied.

\begin{figure}[htbp]\centering
\includegraphics[width=0.94\textwidth]{figures/noise.pdf}
\caption{Synthetic native-gate results with confidence intervals.
The dense large circuit's noisy failure is retained as a feasibility constraint.}
\end{figure}

\section{Reproducibility, measured resource use and storage}
Run all default checks and build the documents with
\begin{quote}\ttfamily\small
docker compose up --build --abort-on-container-exit\\
\hspace*{4mm}--exit-code-from validation
\end{quote}
Use a fresh output directory for each run. The host output location is
controlled by the setting \path{QCENN_ARTIFACTS_DIR}.
The software container executes no QPU submission and does not require
a provider credential. The image build needs package/network access;
the numerical campaigns themselves use local simulators.

The benchmark stage totals in this run are \SuiteWallSeconds{} wall
seconds and \SuiteCPUSeconds{} process CPU seconds. Maximum per-stage
child peak RSS is \SuitePeakMiB{} MiB. Numerical evidence, logs and QPY
exports occupy approximately \ArtifactMiB{} MiB before PDF/source
packaging. These values do not establish EuroHPC node-hour consumption.
Peak RSS is a process measurement, distinct from total node memory and
from the compact simulator's planning estimate.

\begin{center}\small
\begin{tabular}{lrrrr}\toprule
Stage & Result & Wall s & CPU s & Peak child MiB\\\midrule
\input generated/stage_rows.tex 
\bottomrule\end{tabular}
\end{center}

Every campaign preserves its JSON decision fields, raw-count audits where
applicable, circuit exports and log. A JUnit record identifies the unit
test outcome. A SHA-256 manifest detects artifact modification. The
publication source package contains generated table inputs and vector
figures. Software validation and noisy-quality decisions are separate;
the default software exit status does not conceal a hardware-quality
failure because that failure is explicitly reported alongside it.

\section{Work remaining on the allocated emulator}
The minimum next hosting-system run should reproduce this exact source
fingerprint or a new immutable revision and retain the scheduler job ID,
node type/count, allocation account, wall limit, actual wall time, charged
node-hours, memory limits and image digest. Those fields are currently
absent and must not be inferred from a local or CI run. Where Docker is
not permitted, the hosting entity's supported container mechanism or
Python environment should be used with the same locked dependencies.

An illustrative 200-node-hour emulator work plan is shown below, suitable
for the applicant-provided allocation when that amount is confirmed in
the notice. These are planning ceilings, not observed charges or a request
to extend the allocation.

\begin{center}
\begin{tabular}{@{}p{123mm}r@{}}\toprule
Proposed activity & Node-hours\\\midrule
Hosting environment validation and baseline reproduction & 10\\
Extended small-matrix depth/conditioning sweeps & 30\\
Calibration-derived native noise and layout comparisons & 60\\
Structured encoding and circuit-depth reduction experiments & 50\\
Repeated seeds, finite-shot statistics and measurement controls & 30\\
Final regression, artifact verification and technical report & 20\\\midrule
Total proposed ceiling & 200\\\bottomrule
\end{tabular}
\end{center}

The 60-day period must be interpreted from the actual allocation start
notice, which is not present here. Artifact retention should fit the
100-GB allocation by separating curated results from regenerable
intermediates and recording size before extended sweeps. No claim about
remaining time or storage can be made without the hosting account's
actual accounting information.

\section{Proposed first physical experiment and decision gates}
\begin{enumerate}
\item \textbf{Target and calibration.} Obtain the assigned backend name,
native basis, coupling map, calibration snapshot and supported execution
interface. Recompile original source circuits, not a generic-line export.
\item \textbf{Matched emulator.} Execute the compiled order-2/order-4
circuits under the provider calibration-derived noise model; include
state-preparation and reference-verifier controls. Archive all outcomes.
\item \textbf{Bounded QPU pilot.} Start with order 2, then order 4,
using basis and superposition inputs. Record requested and accepted
shots, raw counts, job identifiers, native depth and operations, and
confidence intervals. A preliminary 8192-shot circuit budget may be
used, with subsequent increases driven by accepted-shot uncertainty
and the granted provider allocation.
\item \textbf{Output verification.} Report ancilla acceptance and
reference echo separately. Use signed-state tomography on the small
data registers before claiming a measured full-inverse residual.
No classical residual correction may be presented as quantum output.
\item \textbf{Advance only on evidence.} Apply the declared quality gate
to the controls and inverse circuits. Failures remain in the report.
Order 10 is considered only after small-instance results and device
resource use justify it. Dense orders 100/1000 require a separate
encoding-depth and feasibility assessment.
\end{enumerate}

The requested outcome is an experimentally testable small coherent
inverse-action primitive and a comparison against its calibrated emulator.
It is not a promise of $10^{-12}$ hardware precision, full inverse readout
at dimension 1000, or quantum speedup. Physical access remains subject to
the EuroHPC and hosting-entity assessment.

\section*{Deliverables and disclosure}
The submission package consists of this report, the IEEE-style theoretical
manuscript, source and dependency specifications, complete benchmark
results, unit-test and execution logs, raw-count/circuit evidence and
checksums. Author affiliations and a particular IEEE venue have not been
inferred; those belong to the authors' final submission metadata.

OpenAI ChatGPT/Codex assisted the report and manuscript drafting,
mathematical exposition, workflow implementation and publication scripts.
All plotted measurements are generated programmatically from run outputs.
No AI-generated illustrative image is used as scientific evidence.
The report is prepared for applicant review and is not automatically
sent to EuroHPC or any publication service.

\bibliographystyle{IEEEtran}\bibliography{references}
\end{document}
